\chapter{Úvod}
\pagestyle{plain}
\pagenumbering{arabic}

Každý z nás počas svojho života aspoň raz prezentoval, či to bola prezentácia jednoduchého školského projektu, pracovný pohovor, obhajoba práce alebo profesionálne vystúpenie pred publikom. Výsledok prezentácie má často významný dopad na našu budúcnosť, alebo na zmýšľanie publika. Schopnosť jasne, presvedčivo a sebavedome komunikovať myšlienky je čoraz dôležitejšia nielen pre manažérov či učiteľov, ale aj pre študentov, výskumníkov či pracovníkov v kreatívnych odvetviach.

Mnohí ľudia však pociťujú pri prezentovaní stres, neistotu alebo nedostatok skúseností, čo im bráni v tom, aby naplno využili svoj potenciál. Tradičné formy nácviku, ako sú školenia, kurzy alebo osobné konzultácie, môžu byť časovo náročné, drahé alebo ťažko dostupné. V tomto kontexte vzniká potreba modernejších a flexibilnejších riešení, ktoré by dokázali poskytnúť individuálnu spätnú väzbu a pomohli používateľom zlepšovať sa.

Jedným z takýchto riešení môže byť inteligentný tréner prezentačných schopností, ktorý by využíval moderné technológie, ako sú umelá inteligencia, analýza reči a analýza obrazu na podporu rozvoja komunikačných zručností. Cieľom tejto práce je priblížiť túto tému a ukázať možnosti, ako môže technológia prispieť k efektívnejšiemu a dostupnejšiemu trénovaniu prezentačných schopností.